\begin{exercise}{6}
Show that the category $\mathrm{Set}$ is bicomplete. Hint: General colimits
are constructed as quotients of disjoint unions, and general limits are
constructed as subsets of products.

\begin{proof}
Let $F\colon I\to\mathrm{Set}$ be a small diagram. Form
\[
 Q=\left(\coprod_{i\in\operatorname{Ob}I}F(i)\right)\big/\!\sim,
 \qquad (i,x)\sim(j,F(u)(x))\quad(u\colon i\to j),
\]
where $\sim$ is the equivalence relation generated by these pairs.
The maps $q_i\colon F(i)\to Q$, $x\mapsto[i,x]$, are compatible.
For any compatible family $f_i\colon F(i)\to A$, the formula
$[i,x]\mapsto f_i(x)$ is well-defined, since it respects every
generating pair. It is the unique map $Q\to A$ whose composite with
each $q_i$ is $f_i$. Thus $Q$ is a colimit by Definition~1.1.13.

For the limit, take
\[
 L=\left\{(x_i)\in\prod_{i\in\operatorname{Ob}I}F(i)
       \;\middle|\;F(u)(x_i)=x_j\text{ for every }u\colon i\to j\right\}.
\]
Its projections $p_i\colon L\to F(i)$ are compatible. Any compatible
family $g_i\colon A\to F(i)$ factors uniquely through them via
$a\mapsto(g_i(a))_i$, so $L$ is a limit.
These constructions include $I=\varnothing$: the colimit is the empty
set and the limit is a singleton. Hence $\mathrm{Set}$ has all small
limits and colimits.
\end{proof}
\end{exercise}
